\documentclass[11pt]{article}
\newcommand{\ShortTitle}{INVERSE GALOIS / THE REALIZATION BOUNDARY}
\usepackage[T1]{fontenc}
\usepackage[utf8]{inputenc}
\usepackage{mathpazo}
\usepackage[scaled=0.92]{helvet}
\usepackage{courier}
\usepackage{amsmath,amssymb,mathtools}
\usepackage[a4paper,landscape,left=18mm,right=18mm,top=19mm,bottom=17mm,headheight=14pt]{geometry}
\usepackage{graphicx,microtype,xcolor,fancyhdr,hyperref,booktabs,array}
\definecolor{Ink}{HTML}{243138}
\definecolor{Teal}{HTML}{21676B}
\definecolor{Copper}{HTML}{A45C36}
\definecolor{Soft}{HTML}{E9F0EF}
\definecolor{Gray}{HTML}{627078}
\hypersetup{colorlinks=true,linkcolor=Teal,citecolor=Teal,urlcolor=Teal,pdfauthor={Komorebi P.}}
\urlstyle{same}
\setlength{\parindent}{0pt}
\setlength{\parskip}{0.55em}
\setlength{\emergencystretch}{1.5em}
\linespread{1.04}
\pagestyle{fancy}
\fancyhf{}
\fancyhead[L]{\sffamily\scriptsize\color{Gray}OPERATION KOMOREBI / VISUAL GUIDE}
\fancyhead[R]{\sffamily\scriptsize\color{Gray}\ShortTitle}
\fancyfoot[L]{\sffamily\scriptsize\color{Gray}READER EDITION 1.0 / 9 OCTOBER 2026}
\fancyfoot[R]{\sffamily\small\color{Ink}\thepage}
\renewcommand{\headrulewidth}{0.25pt}
\renewcommand{\footrulewidth}{0pt}
\providecommand{\N}{\mathbb N}
\providecommand{\Z}{\mathbb Z}
\providecommand{\Q}{\mathbb Q}
\providecommand{\R}{\mathbb R}
\providecommand{\C}{\mathbb C}
\providecommand{\eps}{\varepsilon}
\providecommand{\dd}{\mathrm d}
\providecommand{\norm}[1]{\left\lVert #1\right\rVert}
\providecommand{\abs}[1]{\left\lvert #1\right\rvert}
\newcommand{\GuideCover}[4]{%
\hypersetup{pageanchor=false}
\begin{titlepage}\thispagestyle{empty}
{\sffamily\small\color{Teal}OPERATION KOMOREBI\hfill CONJECTURE SERIES #1 / VISUAL GUIDE}\par
\vspace{18mm}
{\fontsize{40}{44}\selectfont\color{Ink}#2}\par
\vspace{4mm}{\LARGE\color{Ink}#3}\par
\vspace{6mm}{\color{Teal}\rule{\linewidth}{1.1pt}}\par
\vspace{4mm}{\large Komorebi P.}\par
{\sffamily\small\color{Gray}READER EDITION 1.0 / 9 OCTOBER 2026}\par
\vfill
\begin{minipage}{0.78\linewidth}\large #4\end{minipage}\par
\vspace{5mm}{\sffamily\small\color{Teal}WORKED EXAMPLES\quad /\quad EXACT COORDINATES\quad /\quad OPEN PROBLEMS}\par
\end{titlepage}\hypersetup{pageanchor=true}}
\newcommand{\GuideSpread}[5]{%
\clearpage
{\fontsize{22}{26}\selectfont\bfseries\color{Ink}#1}\par
\vspace{1mm}{\sffamily\small\color{Gray}#2}\par
\vspace{3mm}
\noindent\begin{minipage}[t]{0.665\linewidth}\vspace{0pt}
\centering\includegraphics[width=\linewidth,height=113mm,keepaspectratio]{#3}
\end{minipage}\hfill
\begin{minipage}[t]{0.295\linewidth}\vspace{0pt}
{\sffamily\small\bfseries\color{Teal}READ THE FIGURE}\par
\small #4
\end{minipage}\par
\vfill
\noindent\colorbox{Soft}{\parbox{\dimexpr\linewidth-2\fboxsep\relax}{\vspace{1.3mm}\small\textbf{Carry forward.} #5\vspace{1.3mm}}}\par
}

\hypersetup{pdftitle={Inverse Galois: A Visual Guide to the Realization Boundary}}
\newcommand{\Gal}{\operatorname{Gal}}
\newcommand{\Tr}{\operatorname{Tr}}
\begin{document}
\GuideCover{09}{Inverse Galois}{A visual guide to the realization boundary}{%
Twenty-two figures follow the passage from a finite group to a Galois extension, from an action to its invariant coordinates, and from a compatible lift to an arithmetic realization. The small examples are complete. The universal existence problem remains open.}

\GuideSpread{01 / What must be realized?}{OPEN QUESTION / one field for every finite group}{figures/01_goal.pdf}{%
The inverse Galois problem over $\Q$ asks whether every finite group $G$ is isomorphic to $\Gal(L/\Q)$ for some finite Galois extension $L$.

Each field at the bottom realizes one specific group. Such constructions teach us how an abstract symmetry becomes an arithmetic object.

The dashed arrow carries the universal quantifier. A family of examples, a faithful representation or a finite computation does not replace that quantifier with a proof.

The monograph develops the examples and then identifies the steps that the operator language still needs.}{%
Keep the base field in view. A realization over an invariant function field and a realization over $\Q$ are different assertions.}

\GuideSpread{02 / The correspondence reverses inclusion}{SOURCE THEOREM / finite Galois correspondence}{figures/02_field_group.pdf}{%
Suppose $L/K$ is finite Galois with group $G$. A subgroup $H\leq G$ fixes the intermediate field $L^H$.

Larger subgroups impose more fixedness and therefore produce smaller fields. The arrows in the two columns point in opposite directions.

The degree formulas keep the dictionary exact:
\[
[L:L^H]=|H|,\qquad
[L^H:K]=[G:H].
\]
The correspondence is a theorem about an extension already given. It does not construct that extension from an arbitrary group. See Milne2022 and the monograph's source contracts.}{%
Subgroup order measures the degree down from $L$. Subgroup index measures the degree up from $K$.}

\GuideSpread{03 / A quadratic group is completely visible}{WRITTEN PROOF / the field $\Q(\sqrt{2})$}{figures/03_quadratic.pdf}{%
The polynomial $x^2-2$ is irreducible over $\Q$ by Eisenstein at $2$. Adjoining one root also supplies its negative, so $L=\Q(\sqrt{2})$ is its splitting field.

There are two automorphisms. The identity fixes the root, while $\sigma$ changes its sign. Applying $\sigma$ twice returns the root.

Thus $\Gal(L/\Q)\cong C_2$. The diagram represents exact conjugation rather than a measured trajectory. The discriminant $8$ is an exact coefficient calculation.}{%
This is a closed realization: a polynomial, its splitting field, its degree and its complete automorphism group.}

\GuideSpread{04 / Independent sign changes make a four-element group}{WRITTEN PROOF / the biquadratic field}{figures/04_biquadratic.pdf}{%
In $L=\Q(\sqrt{2},\sqrt{3})$, the signs of the two square roots can change independently. Neither square root belongs to the quadratic field generated by the other, so $[L:\Q]=4$.

The four sign choices form $C_2\times C_2$. Every nonidentity element has order two.

Its three subgroups of order two fix exactly the three fields shown on the right. Each has degree two over $\Q$.

The square records group composition. It does not assert that every group is a product of sign changes.}{%
A field lattice can be read through its subgroup lattice, once the Galois extension and its action have been established.}

\GuideSpread{05 / An irreducible cubic can already be normal}{WRITTEN PROOF / a cyclic cubic}{figures/05_cyclic_cubic.pdf}{%
Let $\alpha$ be a root of $f(x)=x^3+x^2-2x-1$. The only possible rational roots are $\pm1$, and neither is a root. Thus $f$ is irreducible.

Its discriminant is $49$. An irreducible cubic has a transitive Galois group in $S_3$, and a square discriminant places that group inside $A_3$. The group is therefore $C_3$.

The other roots are $\alpha^2-2$ and $1-\alpha-\alpha^2$. They already lie in $\Q(\alpha)$, and conjugation cycles the three expressions.}{%
The square discriminant certifies an even-permutation group. Irreducibility supplies transitivity. Both facts are needed.}

\GuideSpread{06 / A cubic root is not always the whole splitting field}{EXACT EXAMPLE / approximate root positions}{figures/06_s3_roots.pdf}{%
The root $\sqrt[3]{2}$ generates a cubic field, but its two nonreal conjugates require the cube root of unity $\zeta_3$.

The splitting field is $\Q(\sqrt[3]{2},\zeta_3)$. Its degree is six. Equivalently, $x^3-2$ is irreducible and its discriminant $-108$ is not a rational square, giving group $S_3$.

The dots use approximate complex coordinates. They are a visual aid. The polynomial, degree argument and discriminant establish the Galois group independently of those coordinates.}{%
Always distinguish the field generated by one root from the field containing all roots. A root plot does not certify a Galois group.}

\GuideSpread{07 / Reduction modulo a prime gives a cycle type}{SOURCE INPUT + EXACT FACTORS / the quartic $x^4-x-1$}{figures/07_factor_cycle.pdf}{%
At an unramified prime, the degrees of the irreducible factors record the cycle type of a Frobenius element in the Galois group.

Modulo $2$, the quartic is irreducible. Modulo $7$, it has irreducible factors of degrees $1$ and $3$. Since the discriminant is $-283$, both primes are unramified.

The group contains a four-cycle and a three-cycle. Among the transitive subgroups of $S_4$, this forces $S_4$.

These are cycle types. The finite factors do not number the complex roots or identify specific permutations in a shared root labelling.}{%
Exact finite calculations can support a complete example when a theorem translates their output into the required arithmetic statement.}

\GuideSpread{08 / An abstract group has an exact variable action}{WRITTEN PROOF / the left permutation action}{figures/08_generic_variables.pdf}{%
For a finite group $G$, introduce algebraically independent variables $x_g$, one for each element. Define
\[
h\cdot x_g=x_{hg}.
\]
This gives a faithful action on $F=\Q(x_g:g\in G)$. The product convention is right to left: $h$ acts by left multiplication on each index.

The picture shows the three swaps induced by $(12)$ in $S_3$. It displays exact index changes rather than numerical coordinates.

Faithfulness means that different group elements induce different automorphisms. It still leaves the base field to be identified.}{%
An action supplies symmetry. Its fixed field supplies a base. The two together produce a Galois extension over that base.}

\GuideSpread{09 / The generic construction realizes the group over its fixed field}{SOURCE THEOREM / Artin's theorem with the base stated}{figures/09_invariant_base.pdf}{%
Set $K=F^G$. Artin's fixed-field theorem gives
\[
\Gal(F/K)\cong G,\qquad [F:K]=|G|.
\]
Thus every finite group has a realization over this invariant field.

However, $K$ is not $\Q$. Since $F/K$ is finite, the two fields have the same transcendence degree over $\Q$, namely $|G|$.

Writing only the group isomorphism would hide the issue. The base $F^G$ is part of the statement and part of its proof. See Milne2022.}{%
The generic action closes the group-theoretic construction. It does not close the number-field existence problem.}

\GuideSpread{10 / A route to another base needs its own theorem}{OPEN BOUNDARY / sufficient routes are not equivalences}{figures/10_base_descent.pdf}{%
Rationality of the invariant field is a useful sufficient route. If a suitable invariant base is purely transcendental over $\Q$, specialization can bring the construction back to arithmetic.

Another route starts with a regular realization over $\Q(t)$ and applies Hilbert irreducibility.

Neither route should be read backwards. Failure of a particular rationality claim does not show that the group cannot occur over $\Q$.

The dashed arrow asks for an arithmetic construction with its hypotheses proved. It is not supplied by the action alone.}{%
Noether's rationality problem is one route toward inverse Galois. It is not a necessary condition for every number-field realization.}

\GuideSpread{11 / Averaging over a subgroup makes a projector}{WRITTEN PROOF / exact invariant coordinates}{figures/11_subgroup_projector.pdf}{%
Work in the regular representation $\Q[S_3]$, with basis $\delta_g$ and action $U_h\delta_g=\delta_{hg}$. For $H=\langle(12)\rangle$,
\[
E_H=\frac{1}{|H|}\sum_{h\in H}U_h.
\]
The three pairs are the left $H$-orbits $Hg$. Averaging replaces each pair by its invariant mean.

The image consists of vectors constant on each orbit. Hence $\operatorname{rank}E_H=3$ and $E_H^2=E_H$.

These identities concern a representation already constructed. They are not an arithmetic existence theorem.}{%
In the regular representation, the invariant dimension is the subgroup index $[G:H]$.}

\GuideSpread{12 / Two subgroup projectors need not commute}{EXACT COMPUTATION / written group-algebra argument}{figures/12_s3_matrix.pdf}{%
The basis order is $e,(12),(23),(13),(123),(132)$. A column records the image of its basis vector.

The two averages have rational entries $0$ and $1/2$. Each is a self-adjoint idempotent. Yet their products occur in different orders:
\[
(E_HE_K-E_KE_H)\delta_e
=\frac{\delta_{(123)}-\delta_{(132)}}{4}.
\]
The exact calculation agrees with the written proof. In the group algebra, commuting subgroup averages correspond to the condition $HK=KH$. Here that condition fails.}{%
Projectors can be individually valid while their compositions depend on order. The representation and coefficient domain must remain explicit.}

\GuideSpread{13 / Faithfulness has two different meanings}{WRITTEN PROOF + EXACT EXAMPLE / the standard $S_4$ representation}{figures/13_commutation.pdf}{%
Let $S_4$ permute the four coordinates of its sum-zero subspace $V$. The action of the group on $V$ is faithful.

The four-cycle subgroup $H$ has no nonzero invariant vector in $V$, so $E_H=0$. It commutes with every other operator. Nonetheless, $H$ is not normal, and for $K=\langle(12)\rangle$ one has $HK\ne KH$.

The group-algebra representation loses information even though the group action is faithful. Therefore commuting averages in an arbitrary faithful group representation need not detect permutable subgroups.}{%
The regular representation faithfully sees the group algebra. An arbitrary faithful group representation need not do so.}

\GuideSpread{14 / The regular representation reads the subgroup lattice}{WRITTEN PROOF / six subgroups of $S_3$}{figures/14_coset_lattice.pdf}{%
The nodes are all subgroups of $S_3$. Lines record inclusion, with larger groups above smaller groups.

For every subgroup $H$,
\[
\dim\Q[S_3]^H=[S_3:H].
\]
The identity subgroup fixes the whole six-dimensional space. Each subgroup of order two fixes three dimensions. The subgroup $A_3$ fixes two, and $S_3$ fixes one.

This dimension count follows from left orbit coordinates. If a Galois extension with this group is given, its fixed-field lattice reverses these subgroup inclusions.}{%
The same subgroup index counts regular invariant coordinates and the degree of the associated fixed field over the base. The two interpretations have different inputs.}

\GuideSpread{15 / Trace and norm do not become the same map}{WRITTEN PROOF / a quadratic calculation}{figures/15_trace_norm.pdf}{%
Write $\alpha=a+b\sqrt{2}$ with $a,b\in\Q$. Its conjugate is $a-b\sqrt{2}$.

The subgroup average gives the fixed component $a$. The trace is the sum of conjugates, $2a$. The norm is their product, $a^2-2b^2$.

Averaging and trace are linear. Norm is multiplicative and generally nonlinear as an additive coordinate map.

The shared conjugates explain their relation, but do not identify the three operations. General field constructions must preserve the distinction.}{%
An average can model invariant extraction. A field norm requires the multiplicative structure of the extension.}

\GuideSpread{16 / Regularity makes specialization an arithmetic bridge}{SOURCE THEOREM / Hilbert irreducibility}{figures/16_regular_specialization.pdf}{%
Suppose a finite Galois extension $F/\Q(t)$ has group $G$ and is regular:
\[
F\cap\overline{\Q}=\Q.
\]
Hilbert irreducibility provides suitable rational specializations retaining the group, giving a finite Galois extension of $\Q$ with group $G$.

This is a sufficient implication with a substantial input. Constructing such a regular extension for every finite group is the regular inverse Galois problem.

The figure states the implication. It does not supply the universal regular family. See DebesHarbater2021 and DebesLectureNotes.}{%
A regular realization over $\Q(t)$ implies a realization over $\Q$. The first existence statement still requires proof.}

\GuideSpread{17 / A nonregular toy shows why hypotheses matter}{EXACT EXAMPLES / specialization of $x^3-t$}{figures/17_parameter_bad.pdf}{%
At $t=0$, the polynomial has a repeated root. At $t=1$, its splitting field is $\Q(\zeta_3)$, with group $C_2$. At $t=2$, the splitting field has group $S_3$.

Thus specialization can change both separability and the group.

The generic splitting field of $x^3-t$ contains $\zeta_3$. Its constant field is larger than $\Q$, so this family is not regular over $\Q$.

The toy cannot be inserted into the previous implication as a regular family. Its purpose is to expose that missing hypothesis.}{%
Check constants before invoking regular specialization. A parameterized polynomial is not automatically a regular Galois extension.}

\GuideSpread{18 / A lift must exist in the specified domain}{MODEL + FINITE OBSTRUCTION / embedding problems}{figures/18_embedding_problem.pdf}{%
An embedding problem specifies a profinite group $\Pi$, a continuous epimorphism $\alpha:\Pi\to G$ and a finite-group epimorphism $\pi:E\to G$.

A weak solution is a continuous homomorphism $\beta:\Pi\to E$ satisfying $\pi\beta=\alpha$. A proper solution is also surjective.

For the abstract finite toy $\Pi=G=C_2$, $E=C_4$ and $\alpha=\mathrm{id}$, no lift exists through the parity quotient. An element of order at most two in $C_4$ has trivial parity.

This is not a counterexample involving an absolute Galois group.}{%
Compatibility, continuity and surjectivity are distinct requirements. The actual arithmetic domain cannot be replaced by an arbitrary finite toy.}

\GuideSpread{19 / Splitting gives a weak lift, not automatic properness}{EXACT FINITE TOY / a two-element image in a four-element target}{figures/19_weak_proper.pdf}{%
Take $\Pi=G=C_2$, $E=C_2\times C_2$, and let $\pi$ project onto the second factor. The map
\[
\beta(g)=(0,g)
\]
is a compatible weak lift of the identity.

Its image is the upper row, with two elements. The target has four. Therefore the lift is not proper.

For this finite domain, no proper solution can exist by cardinality. An absolute Galois group is a different domain, so the toy gives no arithmetic impossibility statement. See SchmidtWingberg1998 for the general terminology.}{%
The section of a split extension produces a weak solution. Properness requires enough image to reach the kernel as well.}

\GuideSpread{20 / Identify exactly what the calculus has established}{CLAIM BOUNDARY / group algebra versus arithmetic existence}{figures/20_boundary.pdf}{%
The operator layer gives a precise language for actions, invariant coordinates, subgroup averages and compatibility.

Finite exact computations test these structures in stated examples. They can reveal a false implication and verify a concrete matrix calculation.

The arithmetic layer asks for an actual field over the chosen base, or for a proper solution in an absolute Galois domain. Neither layer is created by declaring an operator identity.

The useful reduction specifies that remaining witness. It must not erase the witness from the final claim.}{%
An algebraic framework can sharpen an obstruction and make a missing implication explicit without resolving the universal existence problem.}

\GuideSpread{21 / The status is part of the mathematical statement}{PROOF / SOURCE / EXACT COMPUTATION / MODEL / OPEN}{figures/21_status_map.pdf}{%
The explicit small fields are complete examples. The generic invariant-base realization is a classical theorem. The projector statements have written derivations and exact finite checks.

The universal regular realization problem over $\Q(t)$ and the inverse Galois problem over $\Q$ remain open.

This publication supplies no compiled Lean certificate for its claims. Exact rational computations are not Lean proofs, and verified examples are not a theorem for every finite group.

Dated preprint context is recorded separately from established inputs in the source contracts.}{%
State a supported conclusion plainly, with its domain and hypotheses. Keep the unresolved universal step visible.}

\GuideSpread{22 / Evidence must support the exact claim being made}{DEPENDENCY MAP / how the publication can be checked}{figures/22_dependency.pdf}{%
Written proofs establish the operator identities and the elementary examples. Primary sources supply the classical theorem inputs with their exact hypotheses.

The reproducible finite computation records polynomial data, modular factors, subgroup matrices and lifting toys. It has a stated finite scope.

The figures carry formulas and examples from those sources. Approximate root coordinates serve only as visual positions.

The remaining universal arithmetic existence question has no supporting proof in this set. The monograph identifies the gap rather than treating the diagram as its solution.}{%
Follow each dependency to its evidence. A finite model may explain the open step, but cannot discharge it.}

\clearpage
{\fontsize{24}{28}\selectfont\bfseries\color{Ink}Reading the guide with the monograph}\par
\vspace{2mm}{\sffamily\small\color{Gray}SOURCE MAP / REPRODUCTION / MATHEMATICAL SCOPE}\par
\vspace{5mm}
\begin{minipage}[t]{.475\linewidth}\vspace{0pt}
{\large\bfseries\color{Teal}A route through the figures}\par
Figures 01--07 develop the problem and five complete small-group constructions. Figures 08--10 explain the generic invariant-base construction and the distinction between its base and $\Q$.

Figures 11--15 give the operator dictionary, including an exact noncommuting pair and a boundary example showing why faithful group actions need not faithfully represent the group algebra.

Figures 16--19 examine the arithmetic inputs: regularity, specialization and the difference between weak and proper lifts. Figures 20--22 summarize the claim boundary and its evidence.

{\large\bfseries\color{Teal}Reproduce the plates}\par
The source bundle contains \texttt{visual/figures.py}, the guide source, all PDF and PNG plates, and an asset hash record in \texttt{figure\_provenance.json}. The exact computation data are stored separately in \texttt{computations/exact\_examples.json}.

Run \texttt{python visual/figures.py} from the project environment. The \texttt{-{}-only} option accepts one or more plate names. The generator checks its independently constructed subgroup matrices against the recorded exact input before rendering.
\end{minipage}\hfill
\begin{minipage}[t]{.475\linewidth}\vspace{0pt}
{\large\bfseries\color{Teal}Primary reading and source contracts}\par
\textbf{Milne2022.} J. S. Milne, \emph{Fields and Galois Theory}. The source contracts locate fixed-field correspondence, finite automorphism groups and the standard polynomial criteria used here.

\textbf{SchmidtWingberg1998.} A. Schmidt and K. Wingberg, \emph{Safarevic's theorem on solvable groups as Galois groups}, arXiv:math/9809211v1. The source distinguishes weak from proper solutions and proves the solvable-group realization theorem.

\textbf{DebesHarbater2021.} P. D\`ebes and D. Harbater, \emph{Pre-Galois theory}, Osaka J. Math. 58 (2021), 71--101. Together with D\`ebes's \emph{Arithm\'etique des rev\^etements de la droite}, it fixes the regularity and Hilbert specialization hypotheses.

\textbf{HuangEtAl2026.} X. Huang et al., \emph{The Mathieu group $M_{23}$ is a Galois group over $\Q$}, arXiv:2608.08538v2. This dated preprint is context, not a dependency of the elementary examples or a proof of the general problem.

The source bundle records exact titles, versions, primary links and theorem locators in \texttt{literature/source\_contracts.md}. The monograph is the place to check the full proof and each source-dependent implication.

{\large\bfseries\color{Teal}What remains open}\par
This edition does not prove that every finite group occurs over $\Q$, or that every finite group has a regular realization over $\Q(t)$. It separates the written operator theorems, complete examples, external inputs and finite models from those universal assertions.
\end{minipage}
\vfill
\colorbox{Soft}{\parbox{\dimexpr\linewidth-2\fboxsep\relax}{\small
\textbf{Edition boundary.} Operation Komorebi, Conjecture Series 09. Author: Komorebi P. Reader edition 1.0, 9 October 2026. The source and evidence bundle preserves the recovered historical inputs alongside this publication.}}
\end{document}
